\chapter{最强大的变革阻力来自群体}


\textit{"群体的智慧是有限的，但群体的力量是巨大的。"}
——古斯塔夫·勒庞


老王我当年在工厂的时候，厂里要推行一个新制度。厂领导开了好几次会，说这制度有多好，能提高多少效率，能节约多少成本。

员工们表面都说"好好好"，会开完了，该怎么样还怎么样。

后来有个老员工跟我说了一句话，我到现在都记得："小王啊，你知道这新制度好，但你不知道这新制度动了谁的蛋糕。"

我问他动了谁的蛋糕。

他说："信息。"

我们那条生产线，以前班长管生产进度，主任管质量检查，厂长管物料调配。你知道这三个"管"是什么意思吗？是"信息"的意思。谁掌握信息，谁就有权力。新制度一来，生产进度AI管了，质量检查AI扫了，物料调配AI算了——这三个"管"的人，突然发现自己的价值没了。

这才是真正的阻力。不是技术难，是利益难。

\section{一、一个组织变革的故事}


2019年，一家全球性的制造企业决定进行一次大规模的数字化转型。这个决定是由CEO和高管团队做出的，他们认为这是公司保持竞争力的必要举措。

转型计划看起来无懈可击：投资了2亿美元，组建了专门的转型团队，引入了最新的AI和自动化技术，制定了详细的实施路线图。

结果：三年后，转型几乎没有任何实质性的进展。2亿美元大部分被消耗在咨询费和试点项目中，核心业务几乎没有受到影响。转型团队被解散，员工士气低落。

事后分析发现，最大的阻力不是来自个人，而是来自一个群体：公司的中层管理者。

这些中层管理者，并不是简单地"抵制变革"。他们抵制的原因，是变革威胁到了他们作为"信息枢纽"的价值。当数字化系统可以直接把信息从一线传到高层时，这些管理者上传下达的角色就变得多余了。

这个故事揭示了一个深刻的真相：\textbf{最强大的变革阻力，往往来自群体，而非个人。}

就像你请了个设计师来家里装修，你跟设计师说"把墙打掉，把厨房改成开放式的"。设计师说"好好好"，然后跟你说了一堆开放式厨房的问题：你做饭油烟大啊、冬天暖气跑掉啊、小孩安全没保障啊。你一听，好像都有道理。但你知道真正的问题是什么吗？设计师不想打墙，因为打墙要重新设计、要改水电图、要算结构——麻烦。但他不直说，他让你觉得打墙本身有问题。

\section{二、为什么群体阻力比个人阻力更强大}


变革管理文献中有一个常见的误解：变革阻力来自那些"不支持变革"的人。如果能够说服这些人，或者替换他们，变革就会成功。

这个假设在个人层面可能是对的。如果关键人物反对变革，换一个人可能解决问题。

但当阻力来自一个群体时，这个方法就失效了。

为什么群体阻力更强大？

\textbf{原因一：群体创造了规范压力}

当一个群体中的大多数人都反对变革时，个体很难独自支持变革——因为那样会让他们成为"异类"。

有个很有意思的规律：同样的技术，放在不同组织里，效果能差好几倍。差距在哪儿？不在技术，在组织。说白了，技术是标尺，组织才是那道坎。

\textbf{原因二：群体放大了恐惧}

个体的恐惧是私密的。当个体的恐惧被放在群体中讨论时，它们被放大、被强化、被合理化。

在上面的案例中，中层管理者们私下聚在一起讨论转型的"风险"，他们的恐惧在彼此的共鸣中不断强化，最终变成了一种集体的信念：变革是危险的。

\textbf{原因三：群体能有效地组织抵制}

分散的个体很难有效地抵制变革。当个体被组织成一个群体时，他们可以协调行动，传播负面的声音，让其他员工也产生疑虑。

\section{三、AI时代的群体阻力特征}


AI转型带来的群体阻力，有一些独特的特征。

\textbf{特征一：AI的威胁更直接}

以往的技术变革，如ERP或CRM的引入，对普通员工来说通常是间接的威胁——它们改变的是工作方式，而不是直接替代人。

AI不同。AI直接威胁到很多人的工作。当员工感受到"AI会取代我"时，这种恐惧是直接的、个人的、难以理性讨论的。

\textbf{特征二：不确定性更大}

AI的能力边界仍在快速演变。即使是专家，也很难准确预测AI在五年后能做什么、不能做什么。

这种不确定性，让AI变革的抵制变得更加复杂——因为员工不只是在抵制已知的威胁，而是在抵制未知的可能性。

\textbf{特征三：焦虑与怀疑交织}

AI变革中，员工的情绪往往是焦虑（AI会不会取代我）和怀疑（AI真的能做到他们说的吗）交织在一起的。

这种混合的情绪，比单纯的恐惧更难处理。

\section{四、如何应对群体阻力}


群体阻力需要用不同于个人阻力的方法来应对。

\textbf{方法一：理解群体，而不是对抗群体}

当面对一个群体的阻力时，最本能的反应是对抗——告诉他们为什么变革是必要的，为什么他们的担忧是错误的。

但这种方法很少有效。当一个群体形成了共同的信念时，外部的理性论证很难动摇它。

更好的方法通常是：理解群体为什么会有这种恐惧，然后在变革设计中回应这种恐惧。

在上面的案例中，如果高管团队在早期就和中层管理者进行深度对话，理解他们真正的担忧（被取代的恐惧），然后在变革设计中给他们新的价值定位（从"信息枢纽"变成"AI输出的判断者"），阻力可能会大大减少。

\textbf{方法二：识别群体中的"关键影响者"}

每个群体中都有一些"关键影响者"——那些不是正式领导但对其他人的态度有重大影响的人。

识别这些关键影响者，理解他们为什么支持或反对变革，往往比直接对整个群体喊话更有效。

\textbf{方法三：创造群体内的分化}

当一个群体看起来完全一致地反对变革时，尝试找到群体内部的裂缝——那些可能支持变革的人，或者那些虽然担忧但不是绝对反对的人。

通过创造群体内的分化，原来的"统一战线"可能被打破。

\section{五、让群体成为变革的推动力}


最理想的情况，不是"克服"群体阻力，而是让群体成为变革的推动力。

这需要把潜在的阻力者变成变革的参与者。

\textbf{让群体参与变革设计}

当员工参与到变革的设计过程中时，他们不只是在"接受"变革，而是在"拥有"变革。

让群体参与变革设计有几个好处：他们的本地知识可以改善变革方案，他们对变革的疑虑可以在设计过程中被回应，他们对变革的承诺可以因为参与而增强。

\textbf{重新定义群体的价值}

当AI改变了一个群体的传统价值时，帮助他们找到新的价值定位，往往比说服他们接受"旧价值已经没用了"更有效。

在上面的案例中，帮助中层管理者看到他们在AI时代的新价值（判断力、人际关系、跨部门协调）而不是告诉他们"你们的信息枢纽角色不重要了"，可能会带来完全不同的结果。

\section{六、结语：变革管理的核心是人和关系}


AI变革，归根结底是人和关系的变革，而非技术变革。

当组织引入AI系统时，它改变的不只是工作流程，还有人与人的关系、人与技术的关系、人与组织的关系。

那些能够理解这一点，并把变革管理的重心放在人和关系上的组织，更可能实现真正的AI转型。

那些只关注技术，而忽视人和关系的组织，即使引入了最先进的AI系统，也很可能只是在表面文章上有所进展。

变革的阻力来自群体，因为群体代表着人性和关系的复杂性。应对这个阻力，需要的不仅是理性的论证，更是对人的理解、尊重和重新定位。

勒庞说"群体的智慧是有限的，但群体的力量是巨大的"——这话在AI变革时代尤其准确。你得承认群体有时候是愚蠢的，但你也得承认，愚蠢的群体可以毁掉任何理性的变革计划。


\textbf{本章小结：}

1. 最强大的变革阻力来自群体：群体创造规范压力，放大的恐惧，有效组织抵制的能力
2. AI时代群体阻力的特征：威胁更直接（直接替代人）、不确定性更大、焦虑与怀疑交织
3. 应对群体阻力的方法：理解群体而非对抗群体、识别关键影响者、创造群体内部分化
4. 让群体成为变革推动力：让群体参与变革设计、重新定义群体的价值定位
5. 变革管理的核心是人和关系——AI变革改变的是人与人的关系、人与技术的关系
